(1) 由 \(\mathbf F=q(\mathbf E+\mathbf v\times\mathbf B)\)，得 \(\dot v_x=qE/m+\omega_c v_y\)、\(\dot v_y=-\omega_c v_x\)，其中 \(\omega_c=qB/m\) 保留电荷符号。横向分量由电场加速，磁场只改变速度方向且瞬时不做功。
(2) 联立线性常微分方程并施加零初速度，得到 \(v_x=(E/B)\sin\omega_ct\)、\(v_y=(E/B)(\cos\omega_ct-1)\)。再积分得 \(x=(E/(B\omega_c))(1-\cos\omega_ct)\)、\(y=(E/(B\omega_c))\sin\omega_ct-(E/B)t\)，公式对带符号 \(\omega_c\) 同时成立。
(3) 在整数个回旋周期 \(2\pi/|\omega_c|\) 上，周期正弦、余弦的平均都为零，故 \(\langle\mathbf v\rangle=-(E/B)\hat{\mathbf y}=(\mathbf E\times\mathbf B)/B^2\)。电荷换号只逆转微观绕圈方向，不改变均匀场的导心漂移速度。
(4) 直接算 \(v^2=2(E/B)^2(1-\cos\omega_ct)\)，于是 \(m v^2/2=m(E/B)^2(1-\cos\omega_ct)\)；同时 \(qEx=(qE^2/(B\omega_c))(1-\cos\omega_ct)\)，由 \(\omega_c=qB/m\) 两式完全相同。\(E\to0\) 时由静止释放始终静止；\(q\) 反号时回旋方向反转，而电场对总动能的净功仍非负。
